\section{Evaluation}
\label{sec:evaluation}

The endpoint in this section is workflow completion time (WCT).  We do not
run a new experiment for this paper revision.  Instead, we re-analyze three
frozen SGLSim result sets produced by an earlier AegisKV prototype.  This
choice recovers an end-to-end metric and preserves unfavorable outcomes, but
creates an important evidence boundary: the WCT results do \emph{not} validate
the current source-owned, independent-tier implementation in
\S\ref{sec:design}.  We ask: (RQ1) does the legacy prototype reduce WCT after
faults; (RQ2) how does it compare with simple backup controls; and (RQ3) do
larger-scale and heavier-failure legacy studies show the same direction?

\subsection{Legacy WCT Methodology}

\begin{table*}[t]
  \centering
  \caption{Frozen historical WCT evidence.  Replotting reads these artifacts
  but launches no simulator, serving runtime, GPU, or scheduler job.}
  \label{tab:legacy-wct-setup}
  \begin{tabular}{@{}p{0.18\textwidth}p{0.37\textwidth}p{0.38\textwidth}@{}}
    \toprule
    Result set & Configuration & Statistical and validity boundary \\
    \midrule
    DP8 policy comparison & Generated shared-prefix workload; 24 workflows,
      651 request records per policy; three soft failures; one seed/trial;
      fixed $\alpha=0.30$; no backup, host-only, peer-GPU, and legacy AegisKV. &
      Direct discrete-event WCT samples with paired workflows.  Small,
      single-trial evidence from a superseded protection policy. \\
    DP64 projection & Qwen3-8B; TP1/PP1/system-DP64; at most 128K context;
      160 workflows and 405 projected retries per policy. & Model projection,
      not measured WCT: CDF shape is quantile-interpolated from 20 DP8
      samples, while a 128K clean run supplies context-time and retry-token
      scales.  It omits correlated outage, co-failure, and queue spillover. \\
    Heavy-failure stress & Qwen3-32B; logical TP1/PP4/system-DP4; 20 workflows;
      offered concurrency 8 or 16; 16 injected failures. & Direct logical-time
      simulator smoke, tagged offline-unvalidated.  It is neither sparse-fault
      evidence nor a physical 16-process deployment; its alpha/beta pressure
      check did not pass. \\
    \bottomrule
  \end{tabular}
\end{table*}

The DP8 comparison is the only frozen data set that places multiple backup
policies on the same paired workload.  For each workflow we report the
empirical WCT CDF and the paired difference from no backup.  A negative
difference favors protection.  The DP64 artifact extrapolates the WCT
distribution and retry-token source at a larger logical deployment.  We
report it only as a projection.  The heavy-failure artifact provides two
direct simulator stress points, but no confidence interval or generalization
beyond its 20 workflows.  We did not combine these heterogeneous result sets,
fit a new model, discard workflows, or select a new operating point after
observing outcomes.  The paper-only renderer binds every source and output by
SHA-256 in \texttt{legacy-wct-submission.manifest.json}.

\subsection{End-to-End WCT Results}

\paragraph{The legacy prototype helps over no backup, but loses to simple
backup controls in the paired DP8 trial.}
Figure~\ref{fig:legacy-wct-policy}a shows the complete 24-workflow WCT
distributions.  Mean WCT is 91.695~s with no backup, 79.734~s with host
backup, 81.731~s with a peer-GPU replica, and 87.756~s with legacy AegisKV.
Thus, legacy AegisKV reduces mean WCT by 3.939~s (4.3\%) and wins 15 of 24
paired workflows, but HostBackup and GPUReplica reduce the mean by 11.961~s
(13.0\%) and 9.964~s (10.9\%), respectively.  The p99 values are 113.733,
99.339, 97.382, and 104.269~s in the same order.  Figure~\ref{fig:legacy-wct-policy}b
also exposes substantial per-workflow dispersion.  This result supports
fault-recovery usefulness relative to recomputation, not baseline dominance.

\begin{figure*}[t]
  \centering
  \includegraphics[width=\textwidth]{figures/legacy-wct-policy-comparison.pdf}
  \caption{Historical direct WCT comparison for one DP8, fixed-$\alpha$
  legacy trial.  (a) Empirical CDF.  (b) Paired policy-minus-no-backup WCT;
  negative is better and every workflow is shown.  The adaptive current
  design was not executed for this artifact.}
  \Description{A WCT CDF compares no backup, host backup, peer GPU replica,
  and legacy AegisKV. A paired-difference panel shows that all three backup
  policies often beat no backup, while host and peer GPU have larger mean
  reductions than legacy AegisKV.}
  \label{fig:legacy-wct-policy}
\end{figure*}

\paragraph{A larger legacy projection predicts an 11.8\% mean-WCT reduction.}
For the projected DP64 Qwen3-8B setting, mean WCT is 4,193.35~s for legacy
AegisKV and 4,753.03~s for no backup, an 11.78\% reduction
(Figure~\ref{fig:legacy-wct-projection}a).  Projected p99 falls from
6,676.35~s to 5,705.16~s (14.55\%).  Of 21.36 million retry tokens, the
projection attributes 26.2\% to legacy L1 reuse, 37.2\% to legacy L2 reuse,
and 36.6\% to recomputation; no backup recomputes all retry tokens
(Figure~\ref{fig:legacy-wct-projection}b).  These values are useful for
capacity planning, but they are generated by the stated extrapolation rather
than 160 observed workflows and therefore are not an empirical scale result.

\begin{figure*}[t]
  \centering
  \includegraphics[width=\textwidth]{figures/legacy-wct-dp64-projection.pdf}
  \caption{Historical \emph{model projection}, not measured WCT, for
  Qwen3-8B at logical DP64 and at most 128K context.  (a) Projected WCT CDF.
  (b) Projected retry-token source.  The projection inherits a 20-sample DP8
  CDF shape and omits correlated failures and shared-queue spillover.}
  \Description{A projected WCT CDF places legacy AegisKV to the left of no
  backup. A stacked bar attributes legacy AegisKV retry tokens to two reuse
  levels and recomputation, while no backup uses only recomputation.}
  \label{fig:legacy-wct-projection}
\end{figure*}

\paragraph{The direction persists in a heavy-failure logical stress smoke.}
At offered concurrency 8, legacy AegisKV reduces simulated mean WCT from
6,259.11 to 5,077.74~s (18.9\%); at concurrency 16, it reduces the mean from
20,794.95 to 14,473.89~s (30.4\%)
(Figure~\ref{fig:legacy-wct-heavyfail}).  The growing gap is consistent with
avoided recomputation becoming more valuable under congestion.  It is not,
however, proof of that causal mechanism: each panel contains only 20
workflows, uses logical PP timing and 16 injected failures, and the artifact's
pressure-validation flag is false.  We therefore classify both panels as
directional stress evidence rather than a formal sensitivity curve.

\begin{figure*}[t]
  \centering
  \includegraphics[width=\textwidth]{figures/legacy-wct-heavyfail-stress.pdf}
  \caption{Historical direct simulator WCT under a heavy-failure logical-PP
  stress smoke (20 workflows per panel).  The apparent larger reduction at
  concurrency 16 is descriptive; these two points are not a calibrated
  scaling law and do not validate the current design.}
  \Description{Two empirical WCT CDFs compare legacy AegisKV and no backup at
  offered concurrency 8 and 16. Legacy AegisKV has lower mean WCT in both
  panels, with a larger descriptive gap at concurrency 16.}
  \label{fig:legacy-wct-heavyfail}
\end{figure*}

\subsection{What the Legacy WCT Data Establish}

Across all three frozen lineages, the earlier AegisKV prototype has lower WCT
than no backup, with mean reductions of 4.3\% in the paired DP8 trial,
11.8\% in the DP64 projection, and 18.9--30.4\% in the heavy-failure smoke.
Only the first and third lineages directly simulate workflow WCT; only the
first compares multiple backup controls; none executes the current design.
Most importantly, the only paired multi-baseline result favors simple host
and peer-GPU backup over legacy AegisKV.  The defensible conclusion is
therefore narrow: preserving reusable KV can reduce end-to-end recomputation
delay, while the incremental WCT value of Slack-aware admission and adaptive
tiering remains unverified for the current implementation.

\subsection{Current-Design Mechanism Telemetry}
\label{sec:current-mechanisms}

We next report a separate evidence track from the frozen current-design
SGLSim campaign.  These are mechanism counters, not WCT and not Estimand~B.
The campaign contains 33 scenarios, 311 policy--scenario cells, ten matched
outer seeds, and 3,110 completed runs over 393 Weka workflows and 98,827
requests.  It spans offered concurrency 32--224; Qwen3-8B TP1/PP1 and
Qwen3-32B TP1/PP4 layouts; 16--128 simulated GPUs; 1--40~GB/s effective
source-to-tier rates; 2--8~GiB Peer-HBM and 16--256~GiB Store-DRAM quotas per
physical GPU/model shard; and 0--41.667 failures per 1,000 physical
GPU-shard-hours.  All values below are bound to lineage
\texttt{4a99c306e8f9}; the paper supplement manifest records that no new run
was launched and that these values make no physical or WCT claim.

\paragraph{Alpha responds to load and available scale.}
Table~\ref{tab:alpha-pressure} reports ten-seed scenario aggregates at the
primary Qwen3-8B configuration.  Raising offered concurrency at 64 simulated
GPUs increases mean foreground pressure from 0.018 to 0.590 and decreases
mean alpha from 0.868 to 0.129.  Holding concurrency at 128 while increasing
the simulated cluster from 16 to 128 GPUs produces the complementary trend:
pressure falls from 0.860 to 0.077 and alpha rises from 0.012 to 0.492.  At
16 GPUs, 27.3\% of steady windows reach alpha zero.  A Qwen3-32B TP1/PP4,
64-GPU stress cell similarly has pressure 0.863, mean alpha 0.019, and 27.6\%
zero-alpha windows.  Thus alpha can collapse to zero under sustained
pressure, rather than being artificially lower-bounded.

\begin{table*}[t]
  \centering
  \caption{Current-design SGLSim Algorithm~3 telemetry.  The concurrency
  sweep fixes 64 GPUs; the cluster sweep fixes $C=128$.  These are mechanism
  aggregates, not WCT or physical measurements.}
  \label{tab:alpha-pressure}
  \begin{tabular}{@{}llrrrrr@{}}
    \toprule
    Sweep & Setting & Pressure & Alpha & Zero $\alpha$ & HBM occ. & DRAM occ. \\
    \midrule
    Concurrency & $C=32$  & 0.018 & 0.868 & 0.0\% & 0.339\% & 0.037\% \\
                & $C=64$  & 0.049 & 0.665 & 0.0\% & 0.666\% & 0.092\% \\
                & $C=128$ & 0.173 & 0.280 & 0.0\% & 0.856\% & 0.253\% \\
                & $C=224$ & 0.590 & 0.129 & 0.0\% & 2.102\% & 0.542\% \\
    \midrule
    Cluster & 16 GPUs  & 0.860 & 0.012 & 27.3\% & 0.115\% & 0.894\% \\
            & 32 GPUs  & 0.411 & 0.089 & 0.0\% & 0.427\% & 0.508\% \\
            & 64 GPUs  & 0.173 & 0.280 & 0.0\% & 0.856\% & 0.253\% \\
            & 128 GPUs & 0.077 & 0.492 & 0.0\% & 0.815\% & 0.126\% \\
    \bottomrule
  \end{tabular}
\end{table*}

Figure~\ref{fig:current-alpha-occupancy} retains one descriptive central-hazard
seed rather than only the aggregates.  Its trajectory contains both
pressure-triggered multiplicative decreases and stable-window additive
increases; $t=0$ is an initialization snapshot, not an AIMD update.  Failure
rate is not an Algorithm~3 input: at $C=128$, changing the declared hazard
from 0 to 41.667 per 1,000 shard-hours leaves mean alpha near 0.28, with only
small differences caused by realized state history.  The occupancy panel also
shows why alpha is not a cache-fill fraction: it is a protection quota
envelope, while actual admitted bytes remain much lower.

\begin{figure*}[t]
  \centering
  \includegraphics[width=\textwidth]{figures/sim-fig05_alpha_occupancy_timeline.pdf}
  \caption{Current-design SGLSim mechanism telemetry.  Top: a descriptive
  lowest-ID central-hazard seed showing alpha, foreground pressure, and
  protection occupancy.  Bottom: rate-zero central full-policy workload reuse
  and intended admission residency.  HBM-only admission is zero in this cell;
  prefix reuse does not deduplicate protection copies.  This figure reports
  neither WCT nor Estimand~B.}
  \Description{Four panels show an alpha and pressure trajectory, Peer-HBM
  and DRAM occupancy, request and token reuse bins, and intended tier
  admission by reuse bin.}
  \label{fig:current-alpha-occupancy}
\end{figure*}

\paragraph{Reuse is abundant, but does not automatically become copy
deduplication.}
Across the trace, 92.63\% of requests have workflow-local prefix reuse and
98.31\% of prompt tokens are reused 64-token hash-ID blocks.  The near-complete
bin contains 52.30\% of requests and 78.46\% of prompt tokens
(Table~\ref{tab:reuse-bins}).  Higher reuse correlates with more DRAM-only
admission in the rate-zero central cell, but the bins are descriptive and the
policy uses Slack, feasibility, capacity, and transfer deadline as well as
reuse.  Therefore this association is not a causal reuse effect.

\begin{table*}[t]
  \centering
  \caption{Current-design rate-zero central reuse detail, pooled over ten
  seeds.  Residency fractions are among intended chunks in each bin; the
  restored-token fraction comes from fixed K64 conditioned probes and is not
  an IID or WCT sample.  HBM-only is 0\% in every row.}
  \label{tab:reuse-bins}
  \begin{tabular}{@{}lrrrrrr@{}}
    \toprule
    Input reuse bin & Requests & Prompt tokens & DRAM-only & HBM+DRAM & Bypass &
      Restored tokens \\
    \midrule
    None          & 7.37\% & 0.39\% & 25.69\% & 74.23\% & 0.08\% & 87.56\% \\
    Partial       & 5.17\% & 1.78\% & 35.69\% & 63.84\% & 0.47\% & 94.93\% \\
    High          & 35.16\% & 19.36\% & 42.22\% & 57.53\% & 0.25\% & 90.17\% \\
    Near-complete & 52.30\% & 78.46\% & 59.41\% & 40.33\% & 0.26\% & 85.78\% \\
    \bottomrule
  \end{tabular}
\end{table*}

\begin{figure*}[t]
  \centering
  \includegraphics[width=\textwidth]{figures/sim-fig06_reuse_and_recovery.pdf}
  \caption{Current-design SGLSim reuse and recovery detail for the rate-zero
  central full policy.  Left: all-request workload composition.  Middle:
  intended residency by reuse bin.  Right: conditioned-probe restored-token
  fraction; strata are descriptive rather than IID.  This is mechanism
  coverage, not a WCT result.}
  \Description{Three panels show request and prompt-token fractions by reuse
  bin, intended DRAM-only and joint HBM-DRAM residency, and restored-token
  fractions for conditioned recovery probes.}
  \label{fig:current-reuse-recovery}
\end{figure*}

\paragraph{The matrix reveals both useful coverage and missing stress.}
At the central point, mean HBM/DRAM occupancy is only 0.856\%/0.253\% and
copy amplification is 1.443$\times$.  Across all 33 full-policy scenarios
(330 seed runs), HBM-only admissions and DRAM evictions are both zero.  Thus,
the campaign exercises independent paths and all outcomes through mechanism
gates, but does not outcome-cover HBM-only retention or DRAM eviction under
tier saturation.  Conditioned restored-token coverage rises from 0\% at
1~GB/s to 87.09\%, 90.92\%, and 93.28\% at 10, 20, and 40~GB/s.  These rates
show transfer-deadline sensitivity, not observed online recovery probability:
they come from fixed live-state probes, copy traffic does not backpressure
foreground inference, and no worker replacement is executed.
